\subsection{FINITE SEQUENCES}

A finite sequence (resp. finite sequence of elements of a set E) is a family (resp. a family of elements of E) whose index set I is a finite set of integers. The number of elements of I is called the length of the sequence.

Let \((t_i)_{i \in I}\) be a finite sequence of length \(n\) . By Proposition 6 of no. 3 there exists a unique isomorphism \(f\) of the interval \([1, n]\) onto the set of integers \(I\) . For each \(k \in [1, n]\) , \(t_{f(k)}\) is said to be the \(k\) th term of the sequence; \(t_{f(1)}\) (resp. \(t_{f(n)}\) ) is the first (resp. last) term of the sequence.

Let \(P\{i\}\) be a relation such that the elements \(i\) for which \(P\{i\}\) is true form a finite set of integers. A finite sequence \((t_i)_{i \in I}\) is then often written \((t_i)_{P\{i\}}\) . For example, when \(I = [a, b]\) , the notation \((t_i)_{a \leqslant i \leqslant b}\) is often used. Under the same conditions, to denote the product of a family of sets \((X_i)_{i \in I}\) the notations

\[
\prod_ {\mathrm{P} \{i \}} \mathrm{X} _ {i} \quad \text { and } \quad \prod_ {i = a} ^ {b} \mathrm{X} _ {i}
\]

are used; and analogous notations for union, intersection, cardinal product, cardinal sum, * composition laws in Algebra *, and so on.

